\chapter{Thế nào là một phép mã hóa SAT tốt?}
\chapterlead{Chất lượng phép mã hóa được đánh giá trên nhiều trục. Đúng đắn là
điều kiện bắt buộc; các trục còn lại gồm kích thước, sức lan truyền, tính địa
phương, chi phí tích hợp và bằng chứng thực nghiệm.}

\index{encoding}\index{soundness}\index{completeness}
\index{unit propagation}\index{arc consistency}\index{channeling}

\flowdiagram{Tính đúng}{Kích thước}{Lan truyền}
  {Cấu trúc bộ giải}{Thực nghiệm}

\section{Phép mã hóa là một quan hệ có giao diện}

Cho ràng buộc \(\mathcal C\) trên tập biến chính \(X\), tập biến phụ \(Y\) và
\[
  F=\operatorname{Enc}(\mathcal C;X,Y).
\]
Tiêu chuẩn cơ bản là tương đương theo phép chiếu:
\[
  F\equiv_X\mathcal C
  \iff
  \forall\alpha_X:\quad
  \alpha_X\models\mathcal C
  \Longleftrightarrow
  \exists\beta_Y:\ \alpha_X\cup\beta_Y\models F.
\]
Hai chiều có tên riêng:
\begin{itemize}
  \item \textbf{soundness (tính đúng)}: mọi mô hình thỏa của \(F\) chiếu thành
        nghiệm của \(\mathcal C\);
  \item \textbf{completeness (tính đầy đủ)}: mọi nghiệm của \(\mathcal C\) mở
        rộng được thành mô hình thỏa của \(F\).
\end{itemize}

Tính đúng thường được chứng minh bằng cách lấy một mô hình \CNF và giải mã.
Tính đầy đủ thường được chứng minh bằng cách xây biến phụ từ một nghiệm miền.
Với bộ đếm, cách xây tự nhiên là đặt \emph{register} (thanh ghi ngưỡng, tức
biến phụ lưu một mức đếm) đúng theo tổng tiền tố thật.

\begin{proposition}[Mẫu chứng minh bằng phép chiếu]
Giả sử bộ giải mã \(D\) ánh mọi mô hình thỏa của \(F\) thành nghiệm của
\(\mathcal C\), và với mọi nghiệm \(s\) của \(\mathcal C\) tồn tại một phép gán biến phụ
\(\beta(s)\) làm \(F\) đúng. Khi đó \(F\equiv_X\mathcal C\).
\end{proposition}

\begin{proof}
Bộ giải mã cho tính đúng. Hàm \(\beta\) cho tính đầy đủ. Hai chiều đúng với
mọi phép gán trên \(X\), nên tương đương theo phép chiếu.
\end{proof}

\section{Hệ phân loại theo cách biểu diễn miền}

Trước khi so hai bộ mã hóa, cần xác định chúng có cùng ``từ vựng Boolean'' hay
không. Với biến miền hữu hạn \(v\in D=\set{d_1,\ldots,d_m}\), bốn họ thường
gặp là:
\begin{description}[style=nextline]
  \item[Direct encoding (mã hóa trực tiếp)]
  \(x_{v,a}\) mang nghĩa \(v=a\). Ràng buộc miền cần đúng-một; xung đột giữa
  hai giá trị được viết trực tiếp.

  \item[Support encoding (mã hóa hỗ trợ)]
  Mỗi lựa chọn \(v=a\) kéo theo một phép tuyển các giá trị hỗ trợ của biến bên
  kia. Cách này thường phản ánh nhất quán cung tốt nhưng các mệnh đề hỗ trợ có
  thể dài.

  \item[Order encoding (mã hóa thứ tự)]
  \(o_{v,t}\) mang nghĩa \(v\ge t\) hoặc \(v\le t\). Bất đẳng thức, khoảng và
  quan hệ trước--sau trở thành phép kéo theo giữa các ngưỡng
  \parencite{tamura2009csp}.

  \item[Log/binary encoding (mã hóa lôgarit/nhị phân)]
  Giá trị dùng \(\lceil\log_2m\rceil\) bit. Số biến nhỏ, nhưng một giá trị hoặc
  khoảng thường không còn là một literal đơn; lan truyền miền có thể đi qua
  mạch sâu hơn.
\end{description}

Mã hóa lai giữ hai góc nhìn, chẳng hạn biến giá trị để chọn một giá trị và biến
thứ tự để biểu diễn một khoảng. Chi phí liên kết hai biểu diễn chỉ đáng bỏ ra
khi góc nhìn thứ hai rút ngắn một họ ràng buộc lớn hoặc tạo suy luận mới. Đây
là kiến trúc được dùng lại trong Chương 7, 9--11.

\begin{center}
\captionof{table}{Các họ biểu diễn miền hữu hạn.}
\label{tab:domain-encodings}
\begin{tabularx}{\textwidth}{@{}>{\bfseries\sffamily}p{.18\textwidth}YYY@{}}
\toprule
Họ & Literal tự nhiên & Phù hợp & Điểm phải kiểm tra\\
\midrule
Trực tiếp & \(v=a\) & Bảng cấm, miền nhỏ & Đúng-một và nhiều cặp xung đột\\
Hỗ trợ & ``\(v=a\) có hỗ trợ'' & Ràng buộc nhị phân dạng bảng & Độ dài mệnh đề hỗ trợ\\
Thứ tự & \(v\ge t\) & Cận, khoảng, hiệu & Miền rộng và liên kết với biến giá trị\\
Lôgarit & Bit của giá trị & Miền lớn, mạch số học & Giá trị ngoài miền, sức lan truyền\\
Lai & Giá trị + ngưỡng/quan hệ & Nhiều kiểu ràng buộc & Chi phí và chiều liên kết\\
\bottomrule
\end{tabularx}
\end{center}

\begin{figure}[tb]
\centering
\begin{tikzpicture}[satfig]
  \node[satlabel,anchor=west] at (0,2.1) {Trực tiếp: một literal cho từng giá trị};
  \foreach \x/\v in {0/1,1.2/2,2.4/3,3.6/4}{
    \node[satbox,minimum width=10mm] at (\x,1.25) {\(v=\v\)};
  }
  \node[satlabel,anchor=west] at (0,.35) {Thứ tự: một chuỗi ngưỡng đơn điệu};
  \foreach \x/\v in {0/2,1.4/3,2.8/4}{
    \node[satstate] (o\v) at (\x,-.55) {\(\ge\v\)};
  }
  \draw[satedge] (o4)--(o3);
  \draw[satedge] (o3)--(o2);
  \draw[decorate,decoration={brace,mirror,amplitude=4pt},Ink]
    (-.5,-1.2)--(3.3,-1.2)
    node[midway,below=5pt,satlabel]{chuyển trạng thái xác định giá trị duy nhất};
\end{tikzpicture}
\caption{Hai từ vựng Boolean cho biến \(v\in\set{1,2,3,4}\).}
\label{fig:direct-order-domain}
\end{figure}

\begin{workedexample}{một bất đẳng thức trong mã hóa trực tiếp và thứ tự}
Cho \(u,v\in\set{0,\ldots,4}\) và \(u+2\le v\). Mã hóa trực tiếp cấm từng cặp
\((a,b)\) với \(a+2>b\):
\[
  \neg x_{u,a}\lor\neg x_{v,b}.
\]
Mã hóa thứ tự dùng \(o_{z,t}\leftrightarrow(z\ge t)\) và các mệnh đề
\[
  \neg o_{u,t}\lor o_{v,t+2},
  \qquad t=1,2.
\]
Nếu một ràng buộc khác suy \(u\ge2\), UP đặt \(o_{u,2}=1\), rồi suy
\(o_{v,4}=1\), tức \(v\ge4\). Một mệnh đề thứ tự đã thay toàn bộ các cặp trực tiếp
có \(u\ge2,v\le3\). Đổi lại, giá trị chính xác của \(v\) chỉ xuất hiện qua điểm chuyển
của chuỗi.
\end{workedexample}

\section{Contract: đặc tả ngữ nghĩa của biến phụ}

Một thanh ghi \(r_{i,j}\) có thể được định nghĩa theo ba mức:
\begin{center}
\captionof{table}{Ba mức đặc tả của literal ngưỡng.}
\label{tab:threshold-contracts}
\begin{tabularx}{\textwidth}{@{}>{\bfseries\sffamily}p{.22\textwidth}YY@{}}
\toprule
Đặc tả & Công thức & Cách sử dụng\\
\midrule
Điều kiện đủ & \(r_{i,j}\rightarrow S_i\ge j\) &
\(r=1\) chứng minh đạt ngưỡng\\
Điều kiện cần & \(S_i\ge j\rightarrow r_{i,j}\) &
\(r=0\) chứng minh chưa đạt ngưỡng\\
Tương đương & \(r_{i,j}\leftrightarrow(S_i\ge j)\) &
Cả hai cực tính đều có ngữ nghĩa\\
\bottomrule
\end{tabularx}
\end{center}

Một \emph{connector} (bộ nối giữa các mô-đun) dùng \(\neg r_{i,j}\) như mệnh
đề ``tổng dưới \(j\)'' phải có chiều cần. Nếu bộ đếm chỉ cung cấp chiều đủ,
bộ nối đó không đúng dù khi đứng riêng hai công thức vẫn có thể
\emph{equisatisfiable} (đồng khả thỏa, tức cùng có nghiệm hoặc cùng vô nghiệm).

\begin{designrule}{Quy tắc ghép mô-đun}
Trước khi nối hai mô-đun, đối chiếu cực tính mà bên sử dụng cần với chiều đặc
tả mà bên cung cấp bảo đảm. Nếu hai phía cần đọc cùng một biến theo hai cách,
phải có mệnh đề liên kết tương ứng.
\end{designrule}

\begin{figure}[tb]
\centering
\begin{tikzpicture}[satfig,node distance=16mm]
  \node[satbox] (p) {Bộ đếm\\cung cấp};
  \node[satstate,right=of p] (r) {\(r_{i,j}\)};
  \node[satbox,right=of r] (c1) {Mô-đun A\\đọc \(r\)};
  \node[satbox,below=11mm of c1] (c2) {Mô-đun B\\đọc \(\neg r\)};
  \draw[satedge] (p)--node[above,satlabel]{``tổng \(\ge j\)''}(r);
  \draw[satedge] (r)--(c1);
  \draw[satedge] (r)--(c2);
  \node[satwarn,rounded corners=2pt,below=11mm of r] (w)
    {Mô-đun B cần chiều\\
     \(\displaystyle S_i<j\Rightarrow\neg r_{i,j}\)};
  \draw[satdash] (w)--(c2);
\end{tikzpicture}
\caption{Đặc tả của literal phụ quyết định cách các mô-đun khác được sử dụng nó.}
\label{fig:provider-client-contract}
\end{figure}

\section{Đo kích thước}

Bốn đại lượng nền tảng là:
\[
  n_v=\card{Y},\qquad
  n_c=\card{F},\qquad
  n_\ell=\sum_{C\in F}\card C,\qquad
  w_{\max}=\max_{C\in F}\card C.
\]
Ngoài độ lớn tiệm cận, cần xem phân bố mệnh đề nhị phân/tam phân, số lần xuất
hiện của biến và bộ nhớ của bộ sinh. Với kỹ thuật hai literal được theo dõi,
tổng số literal thường gần chi phí bộ nhớ hơn số mệnh đề.

Một công thức đóng về \(n_v,n_c\) phải chỉ rõ:
\begin{itemize}
  \item miền tham số, đặc biệt \(k=0,1,n\);
  \item thanh ghi biên nào được bỏ;
  \item khối cuối có thể ngắn;
  \item mệnh đề lôgic nào cần Tseitin hóa thêm.
\end{itemize}

Bộ sinh nên tự xuất thống kê và so với công thức lý thuyết trên các bài toán
nhỏ. Đây là một kiểm thử đơn vị cho phép đếm, đồng thời phát hiện mệnh đề bị lặp hoặc
thanh ghi vô dụng.

\section{Unit propagation: lan truyền đơn vị}

Cho phép gán bộ phận \(A\). \emph{Unit propagation} (lan truyền đơn vị, UP)
lặp việc thỏa một mệnh đề
đơn vị cho
đến điểm cố định hoặc xung đột. Có ba mức phát biểu thường gặp:

\begin{description}[style=nextline]
  \item[UP phát hiện vi phạm]
  Nếu \(A\) không mở rộng được thành nghiệm của \(\mathcal C\), UP trên
  \(F\land A\) tạo xung đột.

  \item[Unit-refutation completeness (URC)]
  Mọi phép gán làm \(F\) không thỏa được đều bị UP chứng minh bất khả thỏa
  trên toàn bộ biến; sách gọi tính chất này là \emph{đầy đủ bác bỏ bằng lan
  truyền đơn vị}.

  \item[Propagation completeness (PC)]
  Mọi literal bị \(F\land A\) kéo theo trong phạm vi quan sát đều được UP suy
  ra, hoặc UP chứng minh bất khả thỏa; đây là \emph{đầy đủ lan truyền}.
\end{description}

Trong một phép mã hóa miền hữu hạn, \emph{generalized arc consistency}
(nhất quán cung tổng quát, GAC) loại mọi giá trị không còn giá trị hỗ trợ.
Phát biểu GAC luôn đi cùng cách Boolean hóa miền và tập literal đại diện cho
giá trị.

Kết quả của Gent và Bacchus cung cấp một cầu nối chính xác: nếu literal Boolean
đại diện trực tiếp cho việc một giá trị còn trong miền, có thể so sánh bao đóng
của lan truyền đơn vị với bao đóng của một bộ lan truyền GAC
\parencite{gent2002arc,bacchus2007gac}. Khi thêm liên kết biểu diễn, đổi từ mã
hóa trực tiếp sang thứ tự hoặc chỉ giữ một chiều Tseitin, tập literal quan sát
thay đổi; phát biểu về GAC vì thế cũng được xác định lại theo tập này.

\begin{example}[AMO từng cặp]
Với
\[
  \operatorname{AMO}(x_1,\ldots,x_n)=
  \bigwedge_{i<j}(\neg x_i\lor\neg x_j),
\]
nếu \(x_i=1\), UP lập tức suy \(\neg x_j\) cho mọi \(j\ne i\). Nếu hai biến
đúng, UP tạo xung đột. AMO từng cặp vì thế có lan truyền rất trực tiếp trên
biến chính, đổi lại \(O(n^2)\) mệnh đề.
\end{example}

\section{Cận dưới kích thước và mặt Pareto}

Yêu cầu lan truyền đặt ra cận dưới cho kích thước phép mã hóa. Với AMO, các
cận dưới đã biết cho thấy một phép mã hóa đầy đủ lan truyền cần số mệnh đề
tuyến tính với hệ số bị chặn, kể cả khi dùng biến phụ
\parencite{kucera2019lower}. Kết quả này dẫn đến ba nguyên tắc:
\begin{enumerate}
  \item xác định rõ lớp mã hóa khi phát biểu về kích thước, biến phụ và sức lan
        truyền;
  \item một cải tiến vài mệnh đề có thể đáng kể nếu nó tiến gần cận dưới;
  \item so sánh kích thước phải giữ cùng tiêu chuẩn lan truyền và cùng ngữ
        nghĩa của đầu ra.
\end{enumerate}

Các kết quả về phép phân rã \textsc{AllDifferent} và ràng buộc đếm toàn cục
còn liên hệ yêu cầu lan truyền với cận biểu diễn của mạch đơn điệu
\parencite{bessiere2009decompositions}. Khi một ràng buộc toàn cục có
bộ lan truyền đạt mức nhất quán mạnh hơn phép mã hóa đa thức đang xét, phép so
sánh tập trung vào phần suy luận được chuyển sang học mệnh đề và chi phí
để đạt phần suy luận đó.

\begin{keyidea}{Mặt Pareto của phép mã hóa}
Một bộ mã hóa được định vị trên bốn trục: kích thước, sức lan truyền, chi phí
xây dựng và khả năng tăng dần. Điểm Pareto là điểm mà cải thiện thêm một trục
sẽ đánh đổi ít nhất một trục khác. \emph{Benchmark} (bộ kiểm chuẩn) đo vị trí
thực nghiệm của từng thiết kế trên mặt Pareto này.
\end{keyidea}

\section{Mức độ phù hợp với bộ giải}

Hai phép mã hóa có cùng tính đúng và kích thước vẫn có thể khác về:
\begin{itemize}
  \item \textbf{tính địa phương}: một ràng buộc cục bộ dùng các biến gần nhau
        hay đi qua một cấu trúc toàn cục;
  \item \textbf{độ sâu kéo theo}: số bước từ quyết định đến xung đột;
  \item \textbf{khả năng học}: xung đột có tạo mệnh đề tái sử dụng được không;
  \item \textbf{đối xứng}: trạng thái phụ có tạo nhiều mô hình tương đương;
  \item \textbf{khả năng tăng dần}: đầu ra có đọc được ở nhiều cận không.
\end{itemize}

Không có một công thức giải tích đơn giản cho các thuộc tính này. Ta dùng
đại lượng thay thế như số mệnh đề nhị phân, đồ thị kéo theo trên bài toán nhỏ,
hoặc số lần lan truyền/xung đột; kết luận cuối cùng cần bộ kiểm chuẩn có kiểm
soát.

\section{Tính mô-đun và mã hóa lại có cấu trúc}

Một mô-đun mã hóa tốt có bốn thành phần:
\[
  \texttt{variables}
  +\texttt{clauses}
  +\texttt{decoder}
  +\texttt{contract}.
\]
Khi nhiều ràng buộc chia sẻ biểu thức con, mã hóa lại có cấu trúc tạo biến phụ
cho phần chung thay vì Tseitin hóa từng công thức riêng
\parencite{haberlandt2023aux}. Lợi ích xuất hiện khi mô-đun thực sự được dùng
chung trong \CNF và bên sử dụng đọc đúng đặc tả.

Ví dụ, nhiều cửa sổ cùng đọc ngưỡng của một tiền tố. Một bộ đếm tiền tố dùng
chung giảm lặp; bộ nối chỉ ghép các ngưỡng ở hai biên. Đây là chủ đề trung tâm
của Chương 5.

\section{Ma trận phát biểu--bằng chứng}

\begin{center}
\captionof{table}{Ma trận phát biểu--bằng chứng cho một phép mã hóa SAT.}
\label{tab:claim-evidence}
\begin{tabularx}{\textwidth}{@{}>{\bfseries\sffamily}p{.24\textwidth}YY@{}}
\toprule
Phát biểu & Bằng chứng lý thuyết & Bằng chứng thực nghiệm\\
\midrule
Mã hóa đúng & Chứng minh hai chiều & Kiểm thử vét cạn/vi sai trên miền nhỏ\\
Mã hóa nhỏ hơn & Công thức đếm & Kích thước \CNF thực tế\\
Lan truyền mạnh hơn & Bổ đề UP/PC/GAC & Bộ đếm lan truyền/xung đột\\
Nhanh hơn & Phân tích cơ chế & Bộ kiểm chuẩn, nhiều \emph{random seed}
(hạt giống ngẫu nhiên), \emph{performance profile} (hồ sơ hiệu năng)\\
Chứng nhận tối ưu & Định lý về cận & Nghiệm và chứng minh ở hai cận kề\\
\bottomrule
\end{tabularx}
\end{center}

Ma trận đưa mỗi ý tưởng vào đúng loại đóng góp: định lý, thiết kế cấu trúc hay
quan sát thực nghiệm.

\section{Một quy trình đánh giá phép mã hóa}

\begin{summarybox}
\begin{enumerate}
  \item Viết rõ ngữ nghĩa của biến chính và biến phụ.
  \item Chứng minh tính đúng và tính đầy đủ theo phép chiếu.
  \item Đếm biến, mệnh đề, literal và độ rộng mệnh đề theo đúng biến thể.
  \item Phân tích một số suy luận cốt lõi của lan truyền đơn vị (UP).
  \item Xây dựng \emph{oracle} (bộ đối chiếu nhỏ trên miền đã biết) để so sánh
        tập nghiệm sau phép chiếu.
  \item Đánh giá trên bộ kiểm chuẩn theo quy trình của Chương 4.
  \item Kiểm tra khả năng kết hợp mô-đun và giải tăng dần.
\end{enumerate}
\end{summarybox}

\section{Bài học thiết kế}

\begin{enumerate}
  \item Tính đúng là quan hệ theo phép chiếu, không phải số mô hình của biến phụ.
  \item Đặc tả ngữ nghĩa quyết định biến phụ có thể được tái sử dụng thế nào.
  \item Đếm literal và cấu trúc mệnh đề quan trọng ngang với số mệnh đề.
  \item Phát biểu về sức lan truyền phải nêu rõ phạm vi literal được quan sát.
  \item Thời gian chạy là kết quả của tương tác giữa phép mã hóa, bộ giải và bài toán thử.
\end{enumerate}

\subsection*{Câu hỏi nghiên cứu}
\begin{enumerate}
  \item Cho một phép mã hóa AMO có biến phụ, viết bộ giải mã và chứng
        minh theo phép chiếu.
  \item Thiết kế một phép thử phân biệt đặc tả điều kiện đủ với đặc tả tương
        đương của thanh ghi ngưỡng.
  \item Chọn ba thước đo để giải thích vì sao một phép mã hóa lớn hơn lại chạy
        nhanh hơn trên các bài toán thử \UNSAT.
\end{enumerate}
